% =====================================================================
%  5.17 MANUAL DE USUARIO
%  Vale un 10% de la nota de Documentacion por si solo (item Doc3).
%  La guia exige TRES partes: requisitos minimos, instalacion y uso.
% =====================================================================

\chapter{Manual de usuario}
\label{cap:manual}

\section{Requisitos mínimos}

La aplicación es una página web que se ejecuta por completo en el
navegador. No requiere cuenta ni instalación de cliente una vez servidos
los ficheros estáticos. El profesional necesita:

\begin{itemize}[leftmargin=1.4em,itemsep=0.2em]
  \item Un navegador con JavaScript activado. La ejecución documentada
  en este capítulo usó Google Chrome 148; no se aporta una matriz de
  compatibilidad certificada con otros navegadores o versiones.
  \item Una pantalla que permita manejar el formulario y el panel de
  resultados. El código incluye adaptaciones de disposición, pero no hay
  campaña de pruebas por resolución.
  \item Conectividad solo para cargar la página y el catálogo de
  criterios. El caso no se envía a ningún servidor.
\end{itemize}

Para instalar o modificar el proyecto se requiere, además, una versión
de Node.js compatible con Angular 20 (\texttt{\^{}20.19.0},
\texttt{\^{}22.12.0} o \texttt{\^{}24.0.0}) y npm. No hace falta base de
datos ni ningún servicio en segundo plano. La ejecución de este manual
usó Node.js~22.14.0 y npm~10.9.7.

\section{Instalación}

Hay dos modos. El de \textbf{producción} es el que acredita este manual:
se genera el artefacto estático y se sirve como ficheros. El de
\textbf{desarrollo} queda para quien vaya a modificar el código.

\subsection{Producción (artefacto estático)}
\label{sec:instalacion-produccion}

Desde la raíz del repositorio, el 1 de septiembre de 2026:

\begin{lstlisting}[numbers=none]
npm install
npx ng build --configuration production
python3 -m http.server 4173 --bind 127.0.0.1 \
  --directory dist/stopp-start-app/browser
\end{lstlisting}

El navegador abre \texttt{http://127.0.0.1:4173}. El resultado de
\texttt{ng build} fue un paquete inicial de 3,53\,MB en
\texttt{dist/stopp-start-app/browser} (advertencias de presupuesto CSS
en los pasos de captura; la compilación terminó con código 0). La
persistencia del caso vive en el \texttt{localStorage} de ese origen;
borrar los datos del sitio equivale a reiniciar.

\subsection{Desarrollo}

\begin{lstlisting}[numbers=none]
npx ng serve --host 127.0.0.1 --port 4200
\end{lstlisting}

El servidor de desarrollo queda en \texttt{http://127.0.0.1:4200} y no
sustituye la instalación de producción. Otras órdenes:
\texttt{npx ng test -{}-watch=false -{}-browsers=ChromeHeadless}
(\cref{anexo:log-karma}) y los comandos del~\cref{anexo:comandos}.

\section{Guía de uso}

El flujo tiene dos pasos y un panel de resultados visible junto a la
captura. No hay pantalla de datos demográficos: se seleccionan
medicamentos, se completan diagnósticos y, si procede, las constantes de
la pestaña «Otros»; los avisos STOPP y START se actualizan al vuelo. Las
figuras de esta sección son capturas de la
ejecución del 1 de septiembre de 2026 sobre el artefacto de producción.

El recorrido habitual es: (1)~marcar los fármacos activos; (2)~leer los
avisos que ya puedan dispararse solo con medicación; (3)~pasar a
diagnósticos; (4)~rellenar la analítica en «Otros» cuando un criterio
dependa de TFGe, presión arterial, iones u otras constantes; (5)~exportar
PDF, copiar el texto, guardar el JSON o reiniciar.

\subsection{Caso ficticio ejecutado el 01/09/2026}
\label{sec:caso-manual}

El siguiente caso no representa a una persona real ni es una recomendación
terapéutica. Se ejecutó el 1 de septiembre de 2026 en Chrome sobre
\texttt{http://127.0.0.1:4173}:

\begin{enumerate}[leftmargin=1.7em,itemsep=0.25em]
  \item Pulsar \emph{Reiniciar} y confirmar \emph{Limpiar todo} para
  partir de un caso vacío.
  \item Mantener vacía la selección de medicamentos.
  \item En \emph{Diagnósticos}, abrir el sistema cardiovascular y el grupo
  \emph{Hipertensión}; marcar \emph{HTA (sin especificar)}.
  \item Comprobar que el panel muestra
  \texttt{START-B1-ANTIHIPERTENSIVO-HTA} (\cref{fig:cap-hta}).
  \item Volver a \emph{Medicamentos}, seleccionar \emph{Amlodipino} y
  comprobar que ese criterio deja de resultar aplicable
  (\cref{fig:cap-amlodipino}).
\end{enumerate}

\paragraph{Salidas observadas.}
\emph{Copiar criterios} (antes de amlodipino) produjo texto con
\texttt{START}, el sistema cardiovascular, el código \texttt{B1} y el
resumen del criterio. \emph{Guardar} descargó un JSON versión~1.0 con
\texttt{diagnoses: ["hta"]} y \texttt{medications} vacío; el fichero
queda en el~\cref{anexo:capturas}. El PDF titulado «Informe STOPP/START»
de esa misma sesión (RF09) se adjunta como fichero
\texttt{memoria/anexos/informe-hta-2026-09-01.pdf}; la
\cref{fig:cap-informe-pdf} del mismo anexo reproduce su primera página.

\begin{figure}[H]
  \centering
  \includegraphics[width=\textwidth]{captura-hta-start-b1.png}
  \caption{Diagnósticos: HTA (sin especificar) marcada y criterio
    START-B1 activado. Ejecución del 01/09/2026 sobre el artefacto de
    producción.}
  \label{fig:cap-hta}
\end{figure}

\subsection{Pantalla de medicamentos}

La~\cref{fig:cap-amlodipino} muestra el primer paso: raíl de sistemas,
grupos propios y relevantes de otros sistemas, y el panel de avisos a la
derecha. Marcar un fármaco foráneo no duplica la selección; \emph{Otro}
añade un ítem con la clase del grupo. Cada pestaña vacía puede marcarse
como revisada.

\begin{figure}[H]
  \centering
  \includegraphics[width=\textwidth]{captura-amlodipino.png}
  \caption{Medicamentos: Amlodipino seleccionado; START-B1 deja de
    aplicar. Ejecución del 01/09/2026 sobre el artefacto de producción.}
  \label{fig:cap-amlodipino}
\end{figure}

\subsection{Pantalla de diagnósticos y analítica}

La~\cref{fig:cap-hta} muestra el segundo paso. Las familias como la
hipertensión son excluyentes. No hay ficha demográfica: la analítica
(TFGe, presión, frecuencia, QTc, iones, TSH) vive en la pestaña
\emph{Otros}; si un valor no se rellena, los criterios que lo leen no
disparan.

\subsection{Resultados y generación del informe}
\label{sec:resultados-manual}

El panel derecho (\cref{fig:cap-resultados}) es el resultado: no hay una
tercera pantalla. Los avisos START van encima; los STOPP, debajo. Cada
bloque muestra un contador y agrupa las tarjetas por sistema. El texto
de cada tarjeta es el \texttt{summary} del criterio; el profesional
juzga si aplican umbrales de dosis o duración.

\begin{figure}[H]
  \centering
  \includegraphics[height=0.26\textheight,keepaspectratio]{captura-resultados.png}
  \caption{Panel de criterios activados del caso HTA, con START-B1 y
    sin avisos STOPP. Detalle de la ejecución del 01/09/2026.}
  \label{fig:cap-resultados}
\end{figure}

Desde cualquier paso se puede exportar PDF, copiar criterios en texto
plano, guardar o cargar JSON v1.0 y reiniciar tras confirmación. Un JSON
ajeno o de versión distinta se rechaza con un mensaje comprensible.

\section{Resolución de errores}

La~\cref{tab:errores-manual} del~\cref{anexo:errores-uso} resume síntomas
de puesta en marcha, exportación y evaluación. No sustituye el
diagnóstico clínico cuando un criterio esperado no aparece.

\section{Preguntas frecuentes}

\paragraph{¿Esta herramienta sustituye al juicio clínico?}
No. Es un cribado: señala criterios STOPP/START que \emph{podrían}
aplicar al caso introducido. La decisión de deprescribir, iniciar o
mantener un tratamiento es siempre del profesional, a la luz de la
situación concreta y de la guía~\cite{omahony2023stopp,sacyl2023stopp}.

\paragraph{¿Dónde se guardan los datos?}
En el \texttt{localStorage} del navegador, en el origen desde el que se
sirve la página (en la ejecución documentada,
\texttt{127.0.0.1:4173}). No hay servidor de aplicación ni cuenta de
usuario. El contenido no sale del equipo salvo que se exporte un JSON o
un PDF.

\paragraph{¿Puedo cargar un caso ya revisado?}
Sí. \emph{Guardar} descarga un JSON v1.0; \emph{Cargar} lo restaura.
El fichero debe ser una exportación de esta aplicación.

